\subsection{Reading the graphs of \(f\) and \(f'\)}

The graph of the derivative records slopes of the original graph. Positive values of \(f'\) indicate increase, negative values indicate decrease, and zeros of \(f'\) indicate horizontal tangents or possible critical points.

\begin{examplebox}[From derivative sign to function shape]
Suppose
\[
f'(x)=(x+1)(x-2).
\]
Then \(f'(x)>0\) for \(x<-1\), \(f'(x)<0\) for \(-1<x<2\), and \(f'(x)>0\) for \(x>2\). Therefore \(f\) increases, then decreases, then increases. The point \(x=-1\) is a local maximum and \(x=2\) is a local minimum.
\end{examplebox}

\begin{interpretationbox}[Reading in the opposite direction]
A rising graph of \(f\) corresponds to \(f'>0\). A horizontal tangent corresponds to \(f'=0\). If the slopes of \(f\) become larger, then \(f'\) is increasing and \(f''>0\). Thus graphs of \(f\), \(f'\), and \(f''\) are three views of the same changing object.
\end{interpretationbox}

\begin{remarkbox}
A zero of \(f'\) is a candidate for an extremum, not an automatic extremum. The sign of \(f'\) on both sides must be checked.
\end{remarkbox}
